\chapter{รากฐานสถาปัตยกรรมเอเจนต์อัจฉริยะและ ReAct Loop}

ระบบปัญญาประดิษฐ์ยุคที่สาม (Third Wave of AI) มีการก้าวกระโดดจากการโต้ตอบข้อความแบบถาม-ตอบธรรมดา สู่การเป็นตัวแทนอัตโนมัติ (Autonomous Agents) ที่มีความสามารถในการวางแผน ใช้เครื่องมือ และตรวจสอบผลลัพธ์ด้วยตนเอง

\section{วิวัฒนาการสู่ Autonomous Multi-Agent Orchestration}

ในระบบเอเจนต์หลายตัว (Multi-Agent System) งานขนาดใหญ่ถูกแบ่งย่อยออกเป็นภารกิจย่อย (Subtasks) และกระจายไปยังเอเจนต์ที่มีความเชี่ยวชาญเฉพาะทาง (Specialized Subagents) โดยมี Supervisor หรือ Orchestrator คอยประสานงานและรวมผลลัพธ์

\section{วงรอบการคิดและการกระทำแบบ ReAct Loop}

แกนกลางของการตัดสินใจของ Agentic AI คือวงรอบ ReAct (Reasoning + Acting) ซึ่งประกอบด้วย 4 ขั้นตอนที่ทำงานต่อเนื่อง:
\begin{enumerate}
    \item \textbf{Thought (การคิด):} การวิเคราะห์สถานะปัจจุบันและเป้าหมายที่ต้องการ
    \item \textbf{Action (การกระทำ):} การเรียกใช้เครื่องมือที่เหมาะสม เช่น การรันคำสั่ง หรือการแก้ไขไฟล์
    \item \textbf{Observation (การสังเกต):} การอ่านผลลัพธ์ที่ได้จากการเรียกใช้เครื่องมือนั้น
    \item \textbf{Reflection (การสะท้อนคิด):} การตรวจสอบว่าผลลัพธ์ที่ได้ถูกต้องตรงตามข้อกำหนดหรือไม่
\end{enumerate}

\begin{theoremlab}{สมการเป้าหมายของการหาเส้นทางคำตอบ (Agent Trajectory Optimization)}
กระบวนการตัดสินใจของเอเจนต์ในขั้นตอนที่ $t$ สามารถแทนได้ด้วยฟังก์ชันความน่าจะเป็น:
\begin{equation}
a_t = \arg\max_{a \in \mathcal{A}} P(a \mid s_t, h_{<t}, \mathcal{G})
\end{equation}
โดยที่ $s_t$ คือสถานะปัจจุบัน, $h_{<t}$ คือประวัติการทำงานในบริบท, และ $\mathcal{G}$ คือเป้าหมายของผู้ใช้
\end{theoremlab}
